% ===========================================================================
%  Capítulo — El lenguaje de dominio específico
% ===========================================================================
\chapter{El Lenguaje de Dominio Espec\'{\i}fico}
\label{cap:dsl}

% ---------------------------------------------------------------------------
\section{Visi\'on general}
\label{sec:vision-general}
% ---------------------------------------------------------------------------

Numerosos problemas cotidianos producen informaci\'on que se organiza
naturalmente en forma de tabla: filas que representan entidades y columnas
que describen sus atributos, relaciones o restricciones.  Los horarios son
un ejemplo paradigm\'atico de esta clase de datos---franjas horarias,
participantes, recursos y conflictos pueden expresarse como tablas
relacionadas---pero la misma estructura aparece en inventarios, resultados
acad\'emicos, planificaci\'on de proyectos y cualquier dominio donde la
informaci\'on admita una organizaci\'on tabular.

El objetivo del lenguaje de dominio espec\'{\i}fico (DSL) desarrollado en
este trabajo es ofrecer un recurso de descripci\'on para esa clase de datos:
un vocabulario declarativo con el que el programador expresa \emph{qu\'e}
contiene cada tabla---columnas, f\'ormulas derivadas, relaciones entre
tablas, reglas de presentaci\'on---sin ocuparse de \emph{c\'omo} se
materializa esa descripci\'on en un formato concreto.

El DSL est\'a implementado como un conjunto de funciones y macros de
Common~Lisp que construyen, en tiempo de carga, un \emph{\'arbol de sintaxis
abstracta} (AST) cuyos nodos son instancias de clases CLOS.  La
canalizaci\'on de compilaci\'on comprende tres etapas claramente separadas:

\begin{enumerate}
  \item \textbf{Expansi\'on de macros.}  Cada forma del DSL
    (p.\,ej.\ \texttt{def-table}, \texttt{exists}) se expande en una llamada
    al constructor CLOS del nodo AST correspondiente.  El resultado es un
    \'arbol cuyos nodos representan las entidades a partir de las cuales el
    sistema es capaz de generar las representaciones visuales de los datos.

  \item \textbf{Generaci\'on de c\'odigo.}  El \'arbol se recorre por el
    backend, que traduce cada nodo a las instrucciones propias del formato de
    salida: expresiones de celda, reglas de presentaci\'on y directivas de
    construcci\'on del documento.  El backend es el \'unico punto que conoce
    los nombres de columna reales, los \'{\i}ndices de fila y los detalles
    de la API de destino.

  \item \textbf{Ejecuci\'on.}  El programa generado se ejecuta, produciendo
    el archivo de salida con los datos, las f\'ormulas derivadas y las reglas
    de presentaci\'on incrustadas.
\end{enumerate}

Un principio central de dise\~no es la \emph{separaci\'on entre intenci\'on
e implementaci\'on}: el DSL expresa la sem\'antica del problema
(``resaltar los profesores que tienen un conflicto de horario'') y el backend
decide el mecanismo concreto seg\'un el formato de destino elegido.  El
programador del DSL no escribe ninguna instrucci\'on espec\'{\i}fica del
formato de salida.

% ---------------------------------------------------------------------------
\section{\'Arbol de sintaxis abstracta}
\label{sec:ast-nodos}
% ---------------------------------------------------------------------------

El \'arbol de sintaxis abstracta (AST) es la representaci\'on intermedia que
separa la capa del DSL de la capa del generador de c\'odigo.  Cada concepto
del problema---una tabla, una expresi\'on de celda, un cuantificador
existencial---queda capturado como un nodo AST con atributos propios.  El
backend opera \'unicamente sobre este \'arbol; no interpreta texto ni
inspecciona las formas del DSL.

Todos los nodos del AST se definen con la macro \texttt{defclass*}, que
genera simult\'aneamente una clase CLOS \texttt{clase-xl-\textit{nombre}} y
una funci\'on constructora \texttt{xl-\textit{nombre}} con par\'ametros por
clave.  Por ejemplo, el nodo del cuantificador existencial se define y
construye de la siguiente manera:

\begin{lstlisting}[language=CommonLisp,
  caption={Patr\'on de definici\'on de nodo AST con defclass*.}]
;; defclass* genera la clase y el constructor xl-expr-exists:
(defclass* xl-expr-exists ()
           (bind-var domain match-keys overlap-cols self-in-cols))

;; El constructor acepta los slots como :keywords:
(xl-expr-exists :bind-var     'r
                :domain       (xl-domain-rows :table-id nil)
                :match-keys   '(dia hora)
                :overlap-cols '(tutor oponente presidente vocal secretario)
                :self-in-cols  nil)
\end{lstlisting}

Los nodos del AST se agrupan en seis categor\'{\i}as seg\'un su papel en
el \'arbol: estructura del workbook, expresiones de celda, navegaci\'on de
fila, cuantificaci\'on existencial, expresiones cross-sheet din\'amicas y
formato condicional con posicionamiento relativo.

% ............................................................................
\subsection{Estructura del workbook}
\label{sec:ast-estructura}
% ............................................................................

Los nodos de estructura representan los contenedores del libro.
Son los nodos de m\'as alto nivel del \'arbol: el workbook agrupa hojas, las
hojas agrupan regiones y las regiones agrupan tablas.  Su jerarqu\'{\i}a
refleja directamente la estructura del documento de salida.

\begin{table}[h!]
\centering
\caption{Nodos de estructura del workbook.}
\label{tab:ast-estructura}
\renewcommand{\arraystretch}{1.3}
\begin{tabular}{p{3.6cm} p{4.8cm} p{5.4cm}}
\toprule
\textbf{Nodo} & \textbf{Atributos} & \textbf{Papel} \\
\midrule
\texttt{xl-workbook}
  & \texttt{name}, \texttt{sheets}
  & Ra\'{\i}z del \'arbol.  \texttt{name} es el nombre del
    archivo de salida; \texttt{sheets} es la lista de hojas. \\
\midrule
\texttt{xl-sheet}
  & \texttt{name}, \texttt{regions},\newline
    \texttt{fixed-expressions}
  & Una hoja del libro.  Contiene regiones horizontales y
    expresiones en posici\'on relativa fuera de las tablas. \\
\midrule
\texttt{xl-region}
  & \texttt{tables}
  & Agrupa tablas colocadas una junto a otra en la misma
    hoja, separadas por una columna en blanco. \\
\midrule
\texttt{xl-table}
  & \texttt{id}, \texttt{col-names},\newline
    \texttt{headers}, \texttt{contenido},\newline
    \texttt{computed}, \texttt{style-rules},\newline
    \texttt{cell-height}, \texttt{cell-width}
  & Plantilla de tabla.  \texttt{col-names} enumera los
    s\'{\i}mbolos DSL de columna; \texttt{computed} asocia
    columnas a expresiones de f\'ormula; \texttt{style-rules}
    almacena las reglas de coloreado condicional. \\
\bottomrule
\end{tabular}
\end{table}

% ............................................................................
\subsection{Expresiones de celda}
\label{sec:ast-expresiones}
% ............................................................................

Las expresiones de celda son los nodos que el backend traduce a
representaciones de f\'ormula en el formato de salida.  Abarcan tanto
referencias a columnas de la misma tabla o de otra (\texttt{xl-expr-column-ref},
\texttt{xl-table-range}) como operaciones de agregaci\'on (\texttt{COUNTIF},
\texttt{COUNTA}, \texttt{SUM}) y literales de texto.  Gracias a esta
representaci\'on, el programador escribe \texttt{(countif rango criterio)}
en lugar de componer manualmente la cadena de f\'ormula con coordenadas
absolutas.

\begin{table}[h!]
\centering
\caption{Nodos de expresi\'on de celda.}
\label{tab:ast-expresiones}
\renewcommand{\arraystretch}{1.3}
\begin{tabular}{p{4.2cm} p{3.8cm} p{5.8cm}}
\toprule
\textbf{Nodo} & \textbf{Atributos} & \textbf{Papel} \\
\midrule
\texttt{xl-expr-column-ref}
  & \texttt{name}, \texttt{context}
  & Referencia a la columna \texttt{name} en la fila actual.
    Se resuelve a la coordenada de esa columna en el formato de salida. \\
\midrule
\texttt{xl-range}
  & \texttt{from-col}, \texttt{to-col}
  & Rango de columnas contiguas de la tabla actual. \\
\midrule
\texttt{xl-table-range}
  & \texttt{table-id},\newline
    \texttt{from-col}, \texttt{to-col}
  & Rango en una tabla espec\'{\i}fica de la regi\'on
    (referencia cross-table).  Lo produce la forma
    \texttt{trange}. \\
\midrule
\texttt{xl-expr-countif}
  & \texttt{count-range},\newline\texttt{criteria}
  & F\'ormula de conteo condicional sobre un rango. \\
\midrule
\texttt{xl-expr-counta}
  & \texttt{count-range}
  & Conteo de celdas no vac\'{\i}as. \\
\midrule
\texttt{xl-expr-sum}
  & \texttt{count-range}
  & Suma sobre un rango. \\
\midrule
\texttt{xl-expr-string}
  & \texttt{value}
  & Literal de texto; produce una cadena fija en la celda. \\
\midrule
\texttt{xl-expr-if}
  & \texttt{test}, \texttt{then}, \texttt{else}
  & Condicional de tres ramas. \\
\midrule
\texttt{xl-expr-equals}
  & \texttt{a}, \texttt{b}
  & Igualdad entre dos expresiones. \\
\midrule
\texttt{xl-expr-and}, \texttt{xl-expr-or}
  & \texttt{a}, \texttt{b}
  & Conjunci\'on y disyunci\'on l\'ogica. \\
\midrule
\texttt{xl-expr-concat}
  & \texttt{a}, \texttt{b}
  & Concatenaci\'on de texto. \\
\midrule
\texttt{xl-expr-table-col-ref}
  & \texttt{table-id}, \texttt{name},\newline\texttt{context}
  & Referencia a una columna de una tabla espec\'{\i}fica de la regi\'on.
    Evita la ambig\"uedad cuando dos tablas comparten el mismo nombre de
    columna. \\
\midrule
\texttt{xl-expr-param-ref}
  & \texttt{name}
  & Referencia a un par\'ametro de instancia almacenado en celda auxiliar. \\
\midrule
\texttt{xl-expr-different}
  & \texttt{a}, \texttt{b}
  & Desigualdad entre dos expresiones. \\
\midrule
\texttt{xl-expr-gt}, \texttt{xl-expr-lt}
  & \texttt{a}, \texttt{b}
  & Comparaci\'on estricta: mayor que y menor que. \\
\midrule
\texttt{xl-expr-gte}, \texttt{xl-expr-lte}
  & \texttt{a}, \texttt{b}
  & Comparaci\'on no estricta: mayor o igual y menor o igual. \\
\midrule
\texttt{xl-expr-add}, \texttt{xl-expr-subtract}
  & \texttt{a}, \texttt{b}
  & Suma y resta aritm\'etica. \\
\midrule
\texttt{xl-expr-multiply}, \texttt{xl-expr-divide}
  & \texttt{a}, \texttt{b}
  & Multiplicaci\'on y divisi\'on aritm\'etica. \\
\midrule
\texttt{xl-expr-time-add}
  & \texttt{a}, \texttt{b}
  & Suma de tiempos como texto:
    \texttt{TEXT(TIMEVALUE(a)+(b/1440),"hh:mm")}. \\
\midrule
\texttt{xl-expr-lookup}
  & \texttt{value-field},\newline\texttt{key-expr}
  & B\'usqueda en la tabla actual: devuelve la columna
    \texttt{value-field} correspondiente a \texttt{key-expr}. \\
\bottomrule
\end{tabular}
\end{table}

% ............................................................................
\subsection{Navegaci\'on de fila y control de visibilidad}
\label{sec:ast-navegacion-fila}
% ............................................................................

Ciertos c\'alculos necesitan referenciar la fila anterior o siguiente de la
misma tabla---por ejemplo, acumular un valor incremental o mostrar una
transici\'on---o bien producir una celda visualmente vac\'{\i}a bajo alguna
condici\'on.  Estos nodos encapsulan esa l\'ogica sin exponer \'{\i}ndices
de fila al programador del DSL.

\begin{table}[h!]
\centering
\caption{Nodos de navegaci\'on de fila y control de visibilidad.}
\label{tab:ast-navegacion}
\renewcommand{\arraystretch}{1.3}
\begin{tabular}{p{3.8cm} p{3.8cm} p{6.2cm}}
\toprule
\textbf{Nodo} & \textbf{Atributos} & \textbf{Papel} \\
\midrule
\texttt{xl-expr-previous-row}
  & \texttt{expr}
  & Eval\'ua \texttt{expr} referenciando la columna en la
    fila anterior de la tabla. \\
\midrule
\texttt{xl-expr-next-row}
  & \texttt{expr}
  & Eval\'ua \texttt{expr} referenciando la columna en la
    fila siguiente de la tabla. \\
\midrule
\texttt{xl-expr-first-row}
  & ---
  & Produce \texttt{TRUE} cuando la fila en evaluaci\'on es la
    primera fila de datos. \\
\midrule
\texttt{xl-expr-non-empty}
  & \texttt{expr}
  & Produce \texttt{TRUE} cuando \texttt{expr} no es vac\'{\i}a
    ni cero. \\
\midrule
\texttt{xl-expr-show-nothing}
  & ---
  & Produce una cadena vac\'{\i}a: la celda queda visualmente
    vac\'{\i}a. \\
\bottomrule
\end{tabular}
\end{table}

% ............................................................................
\subsection{Cuantificaci\'on existencial}
\label{sec:ast-cuantificacion}
% ............................................................................

La cuantificaci\'on existencial es el mecanismo del DSL para expresar
condiciones de la forma ``existe alguna fila que cumpla ciertos criterios''.
Mantenerla como un nodo AST propio, y no como una expresi\'on booleana
expl\'{\i}cita, permite al backend generar la comprobaci\'on m\'as adecuada
seg\'un el contexto: una f\'ormula SUMPRODUCT vectorial en el backend Excel,
u otra implementaci\'on para un backend diferente.

\begin{table}[h!]
\centering
\caption{Nodos de cuantificaci\'on existencial.}
\label{tab:ast-cuantificacion}
\renewcommand{\arraystretch}{1.3}
\begin{tabular}{p{3.6cm} p{4.8cm} p{5.4cm}}
\toprule
\textbf{Nodo} & \textbf{Atributos} & \textbf{Papel} \\
\midrule
\texttt{xl-expr-exists}
  & \texttt{bind-var}, \texttt{domain},\newline
    \texttt{match-keys},\newline
    \texttt{overlap-cols},\newline
    \texttt{self-in-cols}
  & Cuantificador existencial estructurado.  Declara que debe
    existir una fila del dominio que cumpla las condiciones
    expresadas por los slots sem\'anticos.  El backend decide
    el mecanismo de comprobaci\'on seg\'un el tipo de dominio. \\
\midrule
\texttt{xl-domain-rows}
  & \texttt{table-id}
  & Dominio de filas.  \texttt{table-id = nil} indica la
    tabla actual; un s\'{\i}mbolo identifica una tabla
    espec\'{\i}fica de la regi\'on (dominio cross-table). \\
\bottomrule
\end{tabular}
\end{table}

% ............................................................................
\subsection{Expresiones cross-sheet din\'amicas}
\label{sec:ast-cross-sheet}
% ............................................................................

Cuando un libro contiene m\'ultiples hojas con estructura homog\'enea---por
ejemplo, una hoja por grupo o aula---es necesario agregar valores a trav\'es
de ellas.  A diferencia de \texttt{xl-expr-cross-sheet-ref}, que referencia
una hoja con nombre fijo, los nodos de esta secci\'on trabajan con una
\emph{variable de hoja} resuelta en tiempo de compilaci\'on mediante el
entorno din\'amico \texttt{*sheet-env*}, declarado por \texttt{collect-over}.

\begin{table}[h!]
\centering
\caption{Nodos de expresi\'on cross-sheet din\'amica.}
\label{tab:ast-cross-sheet}
\renewcommand{\arraystretch}{1.3}
\begin{tabular}{p{3.8cm} p{4.4cm} p{5.6cm}}
\toprule
\textbf{Nodo} & \textbf{Atributos} & \textbf{Papel} \\
\midrule
\texttt{xl-expr-cross-sheet-ref}
  & \texttt{sheet},\newline\texttt{cell-template}
  & Referencia a una celda en otra hoja con nombre fijo.
    \texttt{cell-template} puede incluir \texttt{\$\{row-num\}} para
    fila din\'amica. \\
\midrule
\texttt{xl-expr-cross-cell}
  & \texttt{sheet}, \texttt{xcol},\newline\texttt{row}
  & Referencia a una celda en la hoja ligada a \texttt{sheet} en
    \texttt{*sheet-env*}. \\
\midrule
\texttt{xl-expr-source-row}
  & \texttt{table-id},\newline\texttt{offset}
  & Fila de una tabla origen mapeada desde la fila l\'ogica actual.
    \texttt{offset} selecciona la subcelda dentro de una fila compuesta. \\
\midrule
\texttt{xl-expr-sheet-id}
  & \texttt{sheet}
  & Identificador can\'onico de la hoja ligada a \texttt{sheet} en
    \texttt{*sheet-env*}. \\
\midrule
\texttt{xl-expr-collect-over}
  & \texttt{groups},\newline\texttt{sheet-var},\newline\texttt{body}
  & Itera \texttt{body} sobre cada hoja de \texttt{groups}, ligando
    \texttt{sheet-var}, y concatena los resultados como lista separada
    por comas. \\
\bottomrule
\end{tabular}
\end{table}

% ............................................................................
\subsection{Formato condicional y posicionamiento}
\label{sec:ast-formato}
% ............................................................................

Estos nodos permiten expresar el coloreado condicional de celdas y la
colocaci\'on de expresiones en posiciones relativas a las tablas.  La
separaci\'on en nodos propios garantiza que el DSL no exponga coordenadas
absolutas: el backend resuelve las posiciones finales en funci\'on del
tama\~no real de los datos.

\begin{table}[h!]
\centering
\caption{Nodos de formato condicional y posicionamiento relativo.}
\label{tab:ast-formato}
\renewcommand{\arraystretch}{1.3}
\begin{tabular}{p{3.6cm} p{4.8cm} p{5.4cm}}
\toprule
\textbf{Nodo} & \textbf{Atributos} & \textbf{Papel} \\
\midrule
\texttt{xl-style-rule}
  & \texttt{rule-condition},\newline
    \texttt{target-columns}
  & Regla de coloreado condicional.  \texttt{rule-condition}
    es un nodo de expresi\'on booleana; \texttt{target-columns}
    enumera las columnas que se colorean cuando la condici\'on
    es verdadera. \\
\midrule
\texttt{xl-anchor}
  & \texttt{table-id},\newline
    \texttt{col-name},\newline
    \texttt{anchor-type}
  & Punto de referencia dentro de una tabla.
    \texttt{anchor-type} es \texttt{:ultima-fila} o
    \texttt{:primera-fila}. \\
\midrule
\texttt{xl-nav}
  & \texttt{anchor}, \texttt{steps}
  & Posici\'on calculada como un ancla m\'as una secuencia de
    pasos direccionales (arriba, abajo, izquierda, derecha). \\
\midrule
\texttt{xl-sheet-fixed-expr}
  & \texttt{pos}, \texttt{expr}
  & Expresi\'on colocada en una posici\'on relativa de la hoja.
    \texttt{pos} es un nodo \texttt{xl-nav}. \\
\bottomrule
\end{tabular}
\end{table}

% ---------------------------------------------------------------------------
\section{Definici\'on de tablas: \texttt{def-table}}
\label{sec:def-table}
% ---------------------------------------------------------------------------

La forma central del DSL para definir datos tabulares es \texttt{def-table}.
Produce una \emph{plantilla} de tabla reutilizable: una descripci\'on
abstracta de las columnas, sus cabeceras, las f\'ormulas que se calculan a
partir de ellas y la regla de coloreado que se aplica sobre sus filas.  La
plantilla no contiene datos; se instancia con datos concretos al ensamblar
el workbook (secci\'on~\ref{sec:workbook}).

\begin{lstlisting}[language=CommonLisp,
  caption={Sintaxis completa de def-table.}]
(def-table nombre-tabla
  ((col1 "Cabecera1") (col2 "Cabecera2") ...)  ; columnas
  :cell-height N       ; filas por fila logica (default 1)
  :cell-width  N       ; columnas por columna logica (default 1)
  :computed            ; formulas derivadas (opcional)
    ((colN (expresion-dsl)))
  :render              ; regla de coloreado (opcional)
    (conditional-rendering
      :condition (expresion-booleana)
      :target-columns (col1 col2 ...)))
\end{lstlisting}

Cada par \texttt{(col "Cabecera")} declara un nombre de columna DSL
(s\'{\i}mbolo) y el texto de su encabezado visual.  El orden de las columnas
es el orden f\'{\i}sico en el documento de salida.

La clave \texttt{:computed} asocia columnas a expresiones DSL que el
compilador traduce a f\'ormulas.  Los valores iniciales de esas columnas
en los datos de entrada se ignoran: la celda siempre mostrar\'a el resultado
de la f\'ormula calculada.

La clave \texttt{:render} acepta una forma \texttt{conditional-rendering}
que especifica bajo qu\'e condici\'on se colorean determinadas columnas.
El mecanismo concreto lo decide el backend.

El listado siguiente define la tabla de profesores del caso de uso
principal, con columna calculada y regla de coloreado condicional:

\begin{lstlisting}[language=CommonLisp,
  caption={Tabla de profesores con columna computed y regla de coloreado.}]
(def-table profesores-table
  ((nombre "") (grado "") (defensas ""))
  :cell-height 1
  :cell-width  1
  :computed
    ;; Cuenta cuantas veces aparece este profesor en los cinco roles
    ((defensas (countif (trange defensas-table tutor secretario)
                        (col nombre))))
  :render
    (conditional-rendering
      :condition
        (exists (r :rows-of defensas-table)
          :self-in  (tutor oponente presidente vocal secretario)
          :matching (dia hora))
      :target-columns (nombre)))
\end{lstlisting}

La variable \texttt{r} en la forma \texttt{exists} es un alias documental
de la fila iterada, an\'alogo a un alias de tabla en SQL.  La regla de
coloreado declara que la celda \texttt{nombre} debe resaltarse cuando el
profesor participa en alguna defensa con conflicto de horario; el mecanismo
lo decide el backend.

% ---------------------------------------------------------------------------
\section{Expresiones de celda}
\label{sec:expresiones-dsl}
% ---------------------------------------------------------------------------

El DSL ofrece formas simb\'olicas para construir expresiones de celda sin
escribir cadenas de f\'ormula directamente.  El compilador resuelve los
nombres de columna a sus coordenadas reales, los rangos a sus l\'{\i}mites
absolutos y las operaciones a su sintaxis en el formato de salida; el
programador solo declara la intenci\'on.

La tabla~\ref{tab:formas-expresion} recoge las formas disponibles.

\begin{table}[h!]
\centering
\caption{Formas de expresi\'on del DSL y su significado.}
\label{tab:formas-expresion}
\renewcommand{\arraystretch}{1.3}
\begin{tabular}{p{5.5cm} p{8.3cm}}
\toprule
\textbf{Forma DSL} & \textbf{Significado} \\
\midrule
\texttt{(col nombre)}
  & Valor de la columna \texttt{nombre} en la fila actual. \\
\midrule
\texttt{(trange tabla-id desde hasta)}
  & Rango que abarca desde la columna \texttt{desde} hasta la
    columna \texttt{hasta} de \texttt{tabla-id}, para todas
    las filas de datos. \\
\midrule
\texttt{(countif rango criterio)}
  & Conteo condicional sobre un rango. \\
\midrule
\texttt{(counta rango)}
  & N\'umero de celdas no vac\'{\i}as en el rango. \\
\midrule
\texttt{(sum-range rango)}
  & Suma de los valores del rango. \\
\midrule
\texttt{(str "texto")}
  & Literal de texto fijo. \\
\midrule
\texttt{(\_equals a b)}
  & Igualdad entre dos expresiones. \\
\midrule
\texttt{(\_and a b)}, \texttt{(\_or a b)}
  & Conjunci\'on y disyunci\'on. \\
\midrule
\texttt{(nav ancla dir n \ldots)}
  & Posici\'on relativa: ancla m\'as pasos direccionales. \\
\bottomrule
\end{tabular}
\end{table}

La forma \texttt{trange} produce \emph{referencias cross-table}: permite que
una expresi\'on en \texttt{profesores-table} cite columnas de
\texttt{defensas-table} sin que el programador conozca las coordenadas
reales.  El compilador las resuelve en funci\'on del orden de declaraci\'on
de las tablas en la regi\'on.

\begin{lstlisting}[language=CommonLisp,
  caption={Expresiones DSL y sus f\'ormulas generadas (backend Excel).}]
;; (col nombre) en profesores-table -> columna K, fila relativa
(col nombre)          ;; -> K4  (en la fila 4)

;; (trange defensas-table tutor secretario) -> rango absoluto de los cinco roles
(trange defensas-table tutor secretario)   ;; -> $B$4:$F$6

;; La formula COUNTIF resultante
(countif (trange defensas-table tutor secretario) (col nombre))
;; -> COUNTIF($B$4:$F$6, K4)
\end{lstlisting}

% ---------------------------------------------------------------------------
\section{El cuantificador existencial \texttt{exists}}
\label{sec:exists-dsl}
% ---------------------------------------------------------------------------

El cuantificador \texttt{exists} permite declarar condiciones de
\emph{existencia} sobre filas de una tabla.  Preguntas del tipo ``existe
otra defensa con el mismo horario'' o ``aparece este profesor en alguna
defensa con conflicto'' son muy frecuentes en problemas de planificaci\'on,
pero expresarlas mediante disyunciones y bucles expl\'{\i}citos produce
c\'odigo verboso y dif\'{\i}cil de leer.  Con \texttt{exists} el
programador declara la intenci\'on sem\'antica a alto nivel; el backend se
encarga de generar la comprobaci\'on correcta.

El cuantificador existe en dos variantes seg\'un el dominio de b\'usqueda.

% ............................................................................
\subsection{Variante same-table: \texttt{:all-rows}}
\label{sec:exists-all-rows}
% ............................................................................

Esta variante busca otra fila dentro de la \emph{misma} tabla que satisfaga
condiciones de coincidencia de valores y solapamiento en columnas de rol.
Se usa para detectar si una fila del propio conjunto tiene alg\'un conflicto
con otra fila del mismo conjunto; por ejemplo, dos defensas programadas en
el mismo d\'{\i}a y hora con al menos un integrante de tribunal en com\'un.

\begin{lstlisting}[language=CommonLisp,
  caption={Cuantificador exists sobre la tabla actual.}]
(exists (r :all-rows)
  :matching    (dia hora)              ; mismos valores en ambas filas
  :any-overlap (tutor oponente presidente vocal secretario))
                                       ; algun valor compartido en estos slots
\end{lstlisting}

La sem\'antica es: ``existe otra fila de esta tabla (distinta de la actual)
cuya columna \texttt{dia} y cuya columna \texttt{hora} coinciden con los de
la fila actual, y que adem\'as comparte al menos un valor en las columnas
de rol''.

En el backend Excel, esta expresi\'on se traduce a una f\'ormula SUMPRODUCT
que multiplica tres factores: coincidencia en las columnas de agrupaci\'on,
exclusi\'on de la fila actual y solapamiento en al menos una columna de rol.

% ............................................................................
\subsection{Variante cross-table: \texttt{:rows-of}}
\label{sec:exists-rows-of}
% ............................................................................

Esta variante busca en las filas de \emph{otra} tabla, relacionando el
valor de la celda actual con columnas de esa tabla.  Se usa cuando el
elemento evaluado (p.\,ej.\ un profesor) y los datos de planificaci\'on
(p.\,ej.\ las defensas) residen en tablas distintas.

\begin{lstlisting}[language=CommonLisp,
  caption={Cuantificador exists sobre otra tabla (cross-table).}]
(exists (r :rows-of defensas-table)
  :self-in  (tutor oponente presidente vocal secretario)
               ; el valor de la celda actual aparece en alguno de estos slots
  :matching (dia hora))
               ; ese slot tiene mas de una defensa programada
\end{lstlisting}

La sem\'antica es: ``existe alguna fila de \texttt{defensas-table} en la que
el valor de la celda actual (el nombre del profesor) aparece en una de las
cinco columnas de rol, y cuyo par (\texttt{dia}, \texttt{hora}) tiene m\'as
de una defensa programada''.

En ambas variantes el alias de fila es puramente documental: el backend
genera directamente una f\'ormula vectorial y no eval\'ua el s\'{\i}mbolo
como variable de programa.  Su papel es an\'alogo al de un alias de tabla
en una sentencia SQL.

% ---------------------------------------------------------------------------
\section{Formato condicional declarativo}
\label{sec:formato-condicional-dsl}
% ---------------------------------------------------------------------------

El coloreado de celdas es una herramienta habitual en los documentos de
planificaci\'on: resaltar en rojo los profesores con conflicto permite
detectar problemas de un vistazo.  El DSL ofrece la forma
\texttt{conditional-rendering} para declarar ese coloreado sin escribir
f\'ormulas ni llamadas a la API de formato.

\begin{lstlisting}[language=CommonLisp,
  caption={Sintaxis de conditional-rendering.}]
(conditional-rendering
  :condition      (expresion-booleana)
  :target-columns (col1 col2 ...))
\end{lstlisting}

Cuando el nodo condici\'on es una instancia de \texttt{xl-expr-exists}, el
backend del caso de demostraci\'on produce una \texttt{FormulaRule} de Excel
incrustada en el archivo \texttt{.xlsx}, en lugar de colorear celdas
est\'aticamente.  Esta regla persiste en el libro y se reevalua
autom\'aticamente cada vez que el usuario edita los datos, sin necesidad de
volver a ejecutar el DSL.

La decisi\'on de usar \texttt{FormulaRule} en lugar de colorear celdas
directamente es exclusiva del backend: el programador del DSL no indica
qu\'e mecanismo de formato se utilizar\'a.

% ---------------------------------------------------------------------------
\section{Ensamblaje del workbook}
\label{sec:workbook}
% ---------------------------------------------------------------------------

Una vez definidas las plantillas de tabla, el programador ensambla el
workbook con tres macros: \texttt{libro} crea el nodo ra\'{\i}z del AST
con el nombre del archivo de salida; \texttt{hoja} agrupa tablas y
expresiones fijas en una pesta\~na del libro; \texttt{tabla} instancia
una plantilla con datos concretos.

\begin{lstlisting}[language=CommonLisp,
  caption={Ensamblaje del workbook con libro, hoja y tabla.}]
(libro nombre-libro
  :filename "Archivo.xlsx"
  :hojas (list
    (hoja "NombreHoja"
      (tabla plantilla1 :data *DATOS-1*)
      (tabla plantilla2 :data *DATOS-2*)
      :fixed-expressions
        (((nav ...) (expresion))
         ((nav ...) (expresion))))))
\end{lstlisting}

Cuando una \texttt{hoja} recibe m\'as de una \texttt{tabla}, el generador
las coloca en la misma regi\'on horizontal, separadas por una columna en
blanco.  Este dise\~no permite que las referencias cross-table
(\texttt{trange}) se resuelvan con coordenadas absolutas correctas.

La clave \texttt{:fixed-expressions} acepta pares de
\texttt{(posici\'on, expresi\'on)}.  La posici\'on se especifica con
\texttt{nav}; su valor se recalcula seg\'un el tama\~no real de los datos,
de modo que los totales de pie de tabla se desplazan autom\'aticamente si se
a\~naden m\'as filas.

% ---------------------------------------------------------------------------
\section{Posicionamiento relativo: \texttt{nav}}
\label{sec:nav-dsl}
% ---------------------------------------------------------------------------

Las expresiones fijas de una hoja no pueden tener coordenadas absolutas
codificadas, pues esas coordenadas dependen de cu\'antas filas de datos
tenga cada tabla en cada ejecuci\'on.  La forma \texttt{nav} resuelve esto:
calcula la posici\'on de una celda como un \emph{ancla} fija en una tabla
m\'as una secuencia de pasos direccionales.

Las anclas disponibles son \texttt{ultima-fila} y \texttt{primera-fila}.
A partir del ancla se encadenan pasos \texttt{abajo}, \texttt{arriba},
\texttt{izquierda} y \texttt{derecha} con su cantidad:

\begin{lstlisting}[language=CommonLisp,
  caption={Sintaxis de nav y anclas de tabla.}]
;; Una fila debajo de la ultima fila de datos
(nav (ultima-fila tabla-id col-name) abajo 1)

;; Una fila debajo y dos columnas a la derecha
(nav (ultima-fila tabla-id col-name) abajo 1 derecha 2)

;; Una fila encima de la primera fila de datos
(nav (primera-fila tabla-id col-name) arriba 1)
\end{lstlisting}

El compilador resuelve el ancla a la coordenada correcta en tiempo de
generaci\'on y produce una referencia absoluta que siempre apunta al lugar
correcto, independientemente del n\'umero de filas de datos.

% ---------------------------------------------------------------------------
\section{Generaci\'on y ejecuci\'on}
\label{sec:generacion-dsl}
% ---------------------------------------------------------------------------

Las dos \'ultimas formas del DSL cierran el ciclo de compilaci\'on:
\texttt{xl-generate} recorre el \'arbol AST y produce el programa de
construcci\'on del libro; \texttt{xl-run-generated} ejecuta ese programa
y crea el archivo de salida.

\begin{lstlisting}[language=CommonLisp,
  caption={Compilaci\'on del AST y ejecuci\'on.}]
;; Recorre el AST y genera el programa de construccion del libro
(xl-generate horario-tesis "horario-tesis.py")

;; Ejecuta el programa generado y crea el archivo .xlsx
(xl-run-generated "horario-tesis.py")
\end{lstlisting}

% ---------------------------------------------------------------------------
\section{Ejemplo completo: Horario de Defensas de Tesis}
\label{sec:ejemplo-completo}
% ---------------------------------------------------------------------------

El caso de uso principal del DSL es la generaci\'on autom\'atica de un
horario de defensas de tesis con dos tablas relacionadas.  La primera tabla,
\texttt{defensas-table}, registra cada defensa con su tribunal (cinco roles:
tutor, oponente, presidente, vocal y secretario), su slot horario y el local.
La segunda tabla, \texttt{profesores-table}, enumera los profesores con su
grado acad\'emico y el n\'umero de defensas en que participan; sus filas se
colorean en rojo cuando el profesor tiene un conflicto de horario.

El listado siguiente es el programa DSL completo; en \'el se observan todas
las construcciones descritas en el cap\'{\i}tulo.

\begin{lstlisting}[language=CommonLisp,
  caption={Programa DSL completo: Horario de Defensas de Tesis.},
  label={lst:tutorial}]
(load "dsl-directo.lisp")
(load "data-tutorial.lisp")   ; *DATOS-DEFENSAS*, *DATOS-PROFESORES*

;; -- TABLA 1: defensas ------------------------------------------------
;; Nueve columnas: cinco de rol (tutor...secretario) mas dia, hora, local.
;; Las cinco columnas de rol son contiguas para poder referenciarlas como
;; rango con (trange defensas-table tutor secretario).
(def-table defensas-table
  ((estudiante "") (tutor "") (oponente "") (presidente "")
   (vocal "") (secretario "") (dia "") (hora "") (local ""))
  :cell-height 1
  :cell-width 1)

;; -- TABLA 2: profesores ----------------------------------------------
;; La columna "defensas" cuenta apariciones del nombre en los cinco roles.
;; El formato condicional marca en rojo a quienes tienen conflicto.
(def-table profesores-table
  ((nombre "") (grado "") (defensas ""))
  :cell-height 1
  :cell-width  1
  :computed
    ((defensas (countif (trange defensas-table tutor secretario)
                        (col nombre))))
  :render
    (conditional-rendering
      :condition
        (exists (r :rows-of defensas-table)
          :self-in  (tutor oponente presidente vocal secretario)
          :matching (dia hora))
      :target-columns (nombre)))

;; -- WORKBOOK ---------------------------------------------------------
(libro horario-tesis
  :filename "Horario_Tesis.xlsx"
  :hojas (list
    (hoja "Defensas"
      (tabla defensas-table   :data *DATOS-DEFENSAS*)
      (tabla profesores-table :data *DATOS-PROFESORES*)
      :fixed-expressions
        (((nav (ultima-fila defensas-table estudiante) abajo 1)
          (str "Total defensas:"))
         ((nav (ultima-fila defensas-table estudiante) abajo 1 derecha 1)
          (counta (trange defensas-table estudiante)))
         ((nav (ultima-fila profesores-table nombre) abajo 1)
          (str "Total profesores:"))
         ((nav (ultima-fila profesores-table nombre) abajo 1 derecha 2)
          (sum-range (trange profesores-table defensas)))))))

(xl-generate horario-tesis "horario-tesis.py")
(xl-run-generated "horario-tesis.py")
\end{lstlisting}

Los datos de entrada contienen tres defensas y seis profesores.  Dos de las
defensas comparten el slot \texttt{Lunes 10:00} con roles solapados.
El programa genera \texttt{Horario\_Tesis.xlsx} con la siguiente estructura:

\begin{itemize}
  \item Columnas de \texttt{defensas-table} con tres filas de datos.
  \item Columna separadora en blanco de la regi\'on.
  \item Columnas de \texttt{profesores-table} con seis filas de datos
    y la columna \texttt{Defensas} calculada mediante \texttt{COUNTIF}.
  \item Fila de totales debajo de cada tabla con posicionamiento din\'amico.
  \item Regla de formato condicional aplicada a la columna \texttt{nombre}:
    los profesores que participan en al menos una defensa cuyo slot tiene
    m\'as de una defensa programada aparecen resaltados en rojo.
\end{itemize}

La regla de coloreado implementa exactamente la sem\'antica del
\texttt{exists} declarado en el DSL.  Al ser una \texttt{FormulaRule}
incrustada en el archivo, cualquier edici\'on posterior de los datos provoca
una reevaluaci\'on autom\'atica del coloreado, sin necesidad de volver a
ejecutar el DSL.
